\section{Reflection positivity and the Euclidean Hilbert space}
\label{ex-reflection}

Reflection positivity is the mechanism that turns Euclidean correlations into
an inner product. A second ingredient is needed to make positive-time
translation contractive: a bound on correlations as the two halves move apart.
We give the complete argument, keeping this bound explicit.

A word \(w=((x_1,v_1),\ldots,(x_n,v_n))\) records finitely many distinct
insertion points and their labels. Write \(S(w)\) for its correlation and
\(\varnothing\) for the empty word. Fix a coordinate axis and call its
coordinate \(t\). Let \(\mathcal W_+\) consist of words whose points have
strictly positive \(t\), including the empty word. Reflection \(\theta\)
reverses the order of a word, reflects the points across \(t=0\), and
applies the reflected conjugate to each label. Assume \(\theta^2=1\).
For \(v,w\in\mathcal W_+\), the concatenation \((\theta v)w\) has no
coincident points: its two parts lie in opposite open half-spaces.

The inputs used in this section are normalization \(S(\varnothing)=1\),
translation invariance, the reflection relation
\(S(\theta w)=\overline{S(w)}\), permutation symmetry on distinct
points, and the following two conditions:
\begin{align}
 &\sum_{i,j=1}^m\overline{c_i}c_j S((\theta w_i)w_j)\ge0
       \quad(w_i\in\mathcal W_+,\ c_i\in\mathbb C),
       \label{ex-rp}\\
 &x\longmapsto S((\theta v)(w+x))
   \text{ is continuous and bounded on }\{x:t(x)\ge0\}
       \quad(v,w\in\mathcal W_+).\label{ex-cap}
\end{align}
The positivity in \eqref{ex-rp} means that the displayed complex number
is real and nonnegative.

\begin{lemma}[From smeared to point reflection positivity]\label{ex-rp-bumps}
Suppose the correlations are continuous at every configuration of distinct
points. Assume reflection positivity for finite sums of smooth compactly
supported tests on positive-time configurations, with arbitrary fixed label
tuples for the separate summands. Then \eqref{ex-rp} holds.
\end{lemma}
\begin{proof}
Fix positive words \(w_1,\ldots,w_m\) and coefficients
\(c_1,\ldots,c_m\). For each nonempty word choose a small configuration
neighborhood in which its points stay distinct and strictly in the positive
half-space. Place a nonnegative smooth bump of integral one in that
neighborhood, centered at the word's positions, with support radius
\(\epsilon\). Such bumps are obtained by translation and rescaling of
\(\exp(-1/(1-|x|^2))\) on the unit ball, zero outside, followed by
division by its positive integral. The empty word has a zero-dimensional
configuration space; use the constant test one there and counting measure
of total mass one.

For each pair \(i,j\), the concatenation of a reflected configuration
near \(w_i\) with a configuration near \(w_j\) has distinct points.
Within either half this follows from the choice of neighborhood; across
halves it follows from the strict time separation. Thus its correlation
kernel is continuous near the pair of centers, even if some spatial
coordinates agree. Write \(K_{ij}^{\epsilon}\) for its integral against
the two bumps. Nonnegativity and unit integrals give
\[
 |K_{ij}^{\epsilon}-S((\theta w_i)w_j)|
 \le \sup\bigl|K_{ij}(x,y)-K_{ij}(x_i,y_j)\bigr|\longrightarrow0,
\]
where the supremum is over the shrinking product support. The kernels are
bounded on these compact supports, so the integrals are well-defined.
Smeared reflection positivity says
\(\sum_{i,j}\overline{c_i}c_jK_{ij}^{\epsilon}\ge0\).
There are finitely many terms, and the nonnegative real half-line is closed
in \(\mathbb C\). Passing to the limit proves \eqref{ex-rp}.
\end{proof}

This argument needs continuity at all distinct-point configurations.
Continuity only where every coordinate is pairwise different would not
suffice: a word and its reflection share spatial coordinates, including in
the diagonal entries of the reflection matrix.

\begin{lemma}[A sufficient separated-correlation bound]\label{ex-cap-proof}
Suppose the correlations are smooth at distinct points and, for each fixed
finite list of labels, there are \(C\ge0\) and \(A\ge0\) such that
\[
 |S_{\boldsymbol v}(\boldsymbol x)|\le C\rho^{-A}
 \quad\text{whenever }0<\rho\le1\text{ and }
       |x_i-x_j|_\infty\ge\rho\ (i\ne j).
\]
Then \eqref{ex-cap} holds.
\end{lemma}
\begin{proof}
Fix \(v,w\). The distances between points within either word are unchanged
by translation. For a reflected point from \(v\) and a point of \(w+x\),
their time separation is the sum of their original positive times plus
\(t(x)\). It is bounded below by a positive number independent of
\(x\) when \(t(x)\ge0\). There are only finitely many pairs. Choose
\(\rho\le1\) smaller than all their positive lower bounds; if there
are no pairs, take \(\rho=1\). The assumed estimate gives the required
uniform bound. The configuration depends continuously on \(x\) and
always remains in the distinct-point set, so smoothness gives continuity.
\end{proof}

\begin{proposition}[Hilbert space and vacuum]\label{ex-gns}
There is a Hilbert space \(\mathcal H_E\) with vectors \([w]\),
\(w\in\mathcal W_+\), whose linear span is dense and such that
\[
 \langle[v],[w]\rangle=S((\theta v)w).
\]
The vector \(\Omega_E=[\varnothing]\) has norm one.
\end{proposition}
\begin{proof}
On the free complex vector space of finitely supported coefficients on
\(\mathcal W_+\), define
\[
 B(c,d)=\sum_{v,w}\overline{c_v}d_w S((\theta v)w).
\]
Reflection, \(\theta^2=1\), and permutation symmetry show
\(B(d,c)=\overline{B(c,d)}\); positivity follows from \eqref{ex-rp}.
For completeness, Cauchy--Schwarz follows by applying positivity to
\(c+\lambda d\): if \(B(d,d)>0\), minimize the resulting quadratic
expression in \(\lambda\); if \(B(d,d)=0\), varying the magnitude and
phase of \(\lambda\) forces \(B(c,d)=0\). Thus
\(|B(c,d)|^2\le B(c,c)B(d,d)\).
Consequently the null vectors \(\mathcal N=\{c:B(c,c)=0\}\) are
orthogonal to every vector and form a subspace. The quotient by
\(\mathcal N\) is an inner-product space. Complete it and map each
word to the class of its basis vector. This gives the stated inner product
and density. Finally
\(\|\Omega_E\|^2=S(\varnothing)=1\).
\end{proof}

Let \(\tau_s w\) translate every point of \(w\) by \(s\) in the positive
time direction, where \(s\ge0\). Define for a finite formal sum \(c\)
the vector \(T_s c=\sum_w c_w[\tau_s w]\). At this point \(T_s\) is
defined on formal sums; descent to the null quotient will follow from its
norm bound. Translation invariance gives
\begin{equation}\label{ex-time-symmetry}
 \langle T_s c,T_t d\rangle=\langle T_0c,T_{s+t}d\rangle.
\end{equation}
Indeed, reflecting the translation by \(s\) shifts the first word by
\(-s\); translating the whole concatenation by \(s\) gives the right
side. Sesquilinearity extends the identity from words to finite sums.

\begin{lemma}[The contraction estimate]\label{ex-contraction}
For every finite sum \(c\) and every \(s\ge0\),
\(\|T_sc\|\le\|T_0c\|\).
\end{lemma}
\begin{proof}
For a single word, \eqref{ex-time-symmetry} gives
\(\|[\tau_sw]\|^2=S((\theta w)(\tau_{2s}w))\).
Condition~\eqref{ex-cap} bounds this expression uniformly in \(s\).
The triangle inequality therefore bounds \(\|T_sc\|\) uniformly for
every finite sum \(c\). Put
\(M=\sup_{s\ge0}\|T_sc\|<\infty\) and \(a=\|T_0c\|\).
By \eqref{ex-time-symmetry} and Cauchy--Schwarz,
\[
 \|T_sc\|^2=|\langle T_0c,T_{2s}c\rangle|\le aM.
\]
Taking the supremum gives \(M^2\le aM\). Since \(M\ge0\), either
\(M=0\), or division by \(M\) gives \(M\le a\). Both cases prove
the result.
\end{proof}

\begin{theorem}[Euclidean translations]\label{ex-euclidean-semigroup}
Positive time translations extend uniquely to a strongly continuous
semigroup \(E_s\), \(s\ge0\), of self-adjoint contractions on
\(\mathcal H_E\). Spatial translations extend to a strongly continuous
unitary representation \(U_y\), \(y\in\mathbb R^3\), commuting with
\(E_s\). Both leave \(\Omega_E\) fixed.
\end{theorem}
\begin{proof}
The contraction estimate sends null vectors to null vectors. It therefore
defines a contraction on the quotient, and then uniquely on its completion.
On word vectors, \(E_0=1\) and \(E_sE_t=E_{s+t}\); continuity of the
bounded operators extends these identities from their dense span to the
whole Hilbert space. Equation~\eqref{ex-time-symmetry} with one time
equal to zero says that \(E_s\) is symmetric. A bounded symmetric operator
defined everywhere is self-adjoint: its adjoint has the same matrix
elements against all vectors, hence is the same operator.

For a spatial shift \(y\), reflection commutes with the shift. Translation
invariance gives
\(\langle[v+y],[w+y]\rangle=\langle[v],[w]\rangle\).
Thus the shift descends to an isometry on the quotient and its completion.
Its inverse is the shift by \(-y\), so it is unitary. The group law and
commutation with \(E_s\) follow on words from commutation of ordinary
translations, and extend by density. The empty word is unchanged by any
translation, proving the vacuum assertions.

It remains to prove strong continuity. For a word \(w\), the square of
\(\|E_sw-E_tw\|\), with word brackets understood, is the sum of the
two diagonal inner products minus twice the real part of the cross term.
By \eqref{ex-time-symmetry} these are values of the continuous function
\(u\mapsto S((\theta w)(\tau_uw))\) at \(2s,2t,s+t\).
Their combination tends to zero as \(s\to t\). The identical norm-square
argument and \eqref{ex-cap} give continuity of spatial translations on
word vectors. Finite linear combinations remain continuous.
For an arbitrary \(\psi\), choose a finite word sum \(v\) with
\(\|\psi-v\|<\varepsilon/3\). For either family \(V_z\), its norm
bound gives
\[
 \|V_z\psi-V_{z_0}\psi\|
 \le2\|\psi-v\|+\|V_zv-V_{z_0}v\|.
\]
The last term is smaller than \(\varepsilon/3\) near \(z_0\), proving
strong continuity.
\end{proof}

This theorem establishes Euclidean Hilbert-space dynamics from the displayed
correlation hypotheses. Relativistic reconstruction also requires field
operators, rotations, locality, boundary distributions and a selected phase
with a unique invariant vacuum. These do not follow merely by naming the
Hilbert space. One possible literature route is the corrected reconstruction
theorem of Osterwalder--Schrader~\cite[\S IV.1]{OS75}; applying that version
requires checking its linear growth condition as well as its covariance,
positivity, symmetry and clustering assumptions. That comparison is still
outstanding here, particularly the passage from bounds for each fixed
label tuple to the required growth control and the treatment of the selected
phase.
